\scanpage{8}

% slanted
\margin{après d)}
\textcircled{e} Pour tout $A\in\mathrm{ob}\,\unsure{\mathscr A|\mathcal A}$, on suppose qu'il y a \inserted{un entier $r=\dim(A)$} et une forme $r$-fois \uline{linéaire}
\[
(\xi_1,\deleted{\unsure{\xi|\cdots}},\xi_r)\leadsto
\deleted{\unsure{\varphi}}I(\xi_1,\cdots,\xi_r)\in\deleted{\mathbb Q}\mathbb Z
\]

2) On suppose $I$ \deleted{\unsure{pos}} str. pos. sur $N^+(A)$, i.e. $\xi_1,\cdots,\xi_r\in N^+(A)\Rightarrow I(\xi_1,\cdots,\xi_r)>0$.

% slanted
\margin{plus \unsure{généralement} si $u:A\longrightarrow B$ est un \unsure{morphisme} \deleted{\unsure{non}} \deleted{\unsure{tel que} $\unsure{u(A)\ne0}$ \illegible{3} \unsure{choisissant} \illegible{2} $\unsure{1\leqslant s\leqslant\dim A}$} \deleted{\unsure{de dim}} et $\eta_1\in N^{\unsure{>|+}}(B)$, alors \unsure{posant|pour} $\xi_1=u^{\unsure{-1|*}}(\eta_1)$, on a $I(\xi_1,\xi_2,\cdots\xi_r)>0$ pour $\xi_2\cdots\xi_r\in N^>(A)$}

3) On suppose que pour tout \unsure{morphisme} $u:A\longrightarrow B$ \deleted{qui est un $\mathbb Q$-isomorphisme, ou comme nous dirons, une \uline{isogénie}, on a $\unsure{r(A)=}$} avec $\dim(A)=\dim(B)$, \deleted{\unsure{et}} il existe un entier \deleted{\unsure{de}} $\unsure{\nu|v}(u)$ tel que
\[
\deleted{\unsure{\xi}} I(u^*(\xi_1),\cdots,u^*(\xi_r))=\nu(u)\,I(\xi_1,\cdots,\xi_r)
\]
pour \unsure{tt|$\forall$} système $\xi_1,\cdots,\xi_r$ [ce qui détermine $\nu(u)$ de façon unique]\deleted{, \unsure{et} cf. plus \unsure{haut}]} et que ($\dim\deleted{\unsure{\mathrm{Ker}\,u}}\inserted{A=\dim B},\ \nu(u)\ne0$) ssi $u$ est un $\mathbb Q$-isomorphisme (``une \uline{isogénie}'').

Prenant en particulier \deleted{$A=B$, et} un $\xi$ dans $N(\deleted{A}B)$ tel que $\varphi_{\unsure{\xi|\eta}}$ soit un isom et $I(\xi,\cdots,\xi)$ soit $>0$ [\deleted{\unsure{on}} \unsure{supposons|supposera} que 2) \unsure{pour cela existe|pour un certain $\xi$}] on \unsure{trouve}
\[
\nu(\alpha)=\frac{1}{I(\xi,\cdots\xi)}I(\alpha^*(\xi),\cdots,\alpha^*(\xi))
\]
$={}$ \uline{forme de degré} $2r$ \uline{en} $\alpha$ \inserted{à coefficients \unsure{entiers}}. Si $A=B$, on aura donc
\[
\nu(\alpha+n\unsure{1|\mathbb 1}_A)=n^{2r}+\deleted{\unsure{\mathrm{tr}}}\sigma_1(\alpha)n^{2r-1}+\cdots+\sigma_{2r}(\alpha)
\]
et on posera
\[
\sigma_1(\alpha)=\mathrm{tr}\,\alpha
\]
tandis qu'on a
\[
\sigma_{2r}(\alpha)=\nu(\alpha)\ ,\ \text{noté aussi }\det\alpha
\]

% archive sheet number: 7

